\documentclass[10pt,a4paper]{amsart}
\usepackage[margin=19mm]{geometry}
\usepackage{amsmath,amssymb,mathtools}
\usepackage[scr=rsfs,cal=euler]{mathalfa}
\usepackage{xfrac}
\usepackage[unicode,hidelinks,bookmarks=false]{hyperref}
\newcommand{\ie}{i.e.\ }
\newcommand{\cf}{cf.\ }
\newcommand{\resp}{respectively\ }
\newcommand{\Subsection}{\subsection}
\providecommand{\nobreakdash}{\nobreak\hbox{-}}
\setlength{\emergencystretch}{2em}
\multlinegap0pt
\usepackage{bm}
\usepackage{bbm}
\usepackage{dsfont}
\usepackage{xspace}
\usepackage{cancel}
\usepackage{mathtools}
\usepackage{amsthm}
\makeatletter
\@ifundefined{theorem}{\newtheorem{theorem}{Theorem}}{}
\@ifundefined{lemma}{\newtheorem{lemma}[theorem]{Lemma}}{}
\@ifundefined{proposition}{\newtheorem{proposition}[theorem]{Proposition}}{}
\makeatother
\DeclareMathOperator{\TV}{TV}
\DeclareMathOperator{\supp}{supp}
\newcommand{\R}{\mathbb{R}}
\newcommand{\Z}{\mathbb{Z}}
\title[Non-constant ground configurations in the disordered ferroma]{Non-constant ground configurations in the disordered ferromagnet}
\author{Michal Bassan, Shoni Gilboa, Ron Peled}
\date{Source: arXiv:2309.06437v3 (CC BY 4.0)}
\begin{document}
\maketitle

\noindent\emph{Source and licence.} Non-constant ground configurations in the disordered ferromagnet. Version arXiv:2309.06437v3 (CC BY 4.0), distributed under the Creative Commons Attribution 4.0 International licence, as recorded in the arXiv licence metadata for this exact version and shown on the abstract page https://arxiv.org/abs/2309.06437v3. Full rights notice: licenses/arxiv-2309.06437-CC-BY-4.0.txt.\par\smallskip

\noindent\textit{Editorial scope.} Counted unit: the Proposition whose source label is \texttt{prop:functional\_diff} in the article (source0.tex lines 1610--1621, source label \texttt{prop:functional\_diff}), delivered together with the article's own complete original proof of it (source0.tex lines 1705--1761, one \texttt{proof} environment of 57 lines). No mathematics is written, repaired or re-derived: statement and proof are the article's own text line for line. The proof is detached from the statement in the source: the article states the result at lines 1610--1621 and prints its proof at lines 1705--1761, and the proof environment's own header names the statement, so the association is the article's own. Required context retained uncounted: lem:Ron (lines 1643--1648); lem:rough (lines 1650--1657); lem:diff\_rough (lines 1659--1666). Every definition, display and label of those lines is retained unchanged. Omitted from this package and not redistributed: 1--1609, 1622--1642, 1649--1649, 1658--1658, 1667--1704, 1762--4061. Nothing inside a retained excerpt is omitted or altered: every retained body is delivered exactly as the article's lines are written, so no omission receipt of any kind accompanies this package. Citations: the retained excerpts carry no citation command at all, so no bibliography entry of the article is transcribed and no reference is redistributed. Reference adaptation: none was needed and none was made; every reference inside the delivered unit resolves inside it, so no unresolved reference or citation is produced. Printed slips kept literal: no printed word, symbol, hypothesis or number of the counted statement or of its proof was altered. Layout and class: the article's own class is \texttt{amsart} with packages including amsmath, amsfonts, amssymb, amsthm, a4wide, url, mathscinet, mathdots, graphicx, float, hyperref, color, amsrefs, xcolor; the wrapper is the lane's pinned \texttt{amsart} editorial wrapper at 10pt on A4, which restates the theorem-like environments the retained text needs and adds the article's own macro definitions that the retained text needs (\texttt{\textbackslash R}, \texttt{\textbackslash TV}, \texttt{\textbackslash Z}, \texttt{\textbackslash supp}), with extra wrapper packages (bm, bbm, dsfont, xspace, cancel, mathtools). The delivered numbering is the unit's own. Assets: no retained span contains a figure environment, a graphics-inclusion command, a PDF-page inclusion, a table or a TikZ picture; no image file is copied into this package, the package declares no asset dependency and no image file is redistributed with it. Licence: version arXiv:2309.06437v3 of this article is distributed under the Creative Commons Attribution 4.0 International licence, as recorded in the arXiv licence metadata for this exact version and shown on the abstract page https://arxiv.org/abs/2309.06437v3. A full rights notice is delivered at licenses/arxiv-2309.06437-CC-BY-4.0.txt. Characters: every retained excerpt is pure ASCII, so the delivered engine, XeLaTeX with its default Latin Modern text font, reports no missing character.
\begin{lemma}\label{lem:Ron}
For every $0<\alpha<\frac{1}{2}$ and $u\in\Z^d$ such that $\lvert\tau^{\rm rough}_N(N\,u)-\tau_N(N\,u)\rvert>\alpha$, it holds that 
\begin{equation*}
\TV\left(\tau; Q_N(N\,u)\right)\geq \frac{2\alpha}{3} N^{d-1}.
\end{equation*}
\end{lemma}

\begin{lemma}
\label{lem:rough}
For every $\{u,v\}\in E(\Z^d)$, 
\begin{equation*}
\left\lvert\tau^{\rm rough}_N(N\,u)-\tau^{\rm rough}_N(N\,v)\right\rvert
\leq\frac{1}{N^{d-1}}\TV\left(\tau; Q_N(N\,u)\cup Q_N(N\,v)\right).
\end{equation*}
\end{lemma}

\begin{lemma}
\label{lem:diff_rough}
Let $w\in{\mathbb Z}^d$, let $g$ be a function from $B:=Q_2(2\,w)$ to $\R$, and let $\mu:=\frac{1}{2^d}\sum_{u\in B}g(u)$. Then,
$$
\sum_{u\in B} \lvert g(u)-\mu \rvert
\leq\sum_{\{u,v\}\in E(\Z^d)\cap\binom{B}{2}} \lvert g(u)-g(v) \rvert.
$$
\end{lemma}

\begin{proposition}\label{prop:functional_diff}
For every shift $\tau$,
$$
\lvert\supp(\tau-\tau_2)\rvert\leq\lVert \tau-\tau_2\rVert_1
\leq 2\, \TV(\tau) 
$$
and moreover, for every positive integer $N$,
$$
\lvert\supp(\tau_N-\tau_{2N})\rvert\leq\lVert \tau_N-\tau_{2N}\rVert_1
\leq (4d+9)N\, \TV(\tau). 
$$
\end{proposition}

\begin{proof}[Proof of Proposition \ref{prop:functional_diff}]

It is clear that $|\supp(\tilde{\tau})|\le \|\tilde{\tau}\|_1$ for every shift $\tilde{\tau}$, so that we only need to bound the $\ell_1$ norms appearing in the statement.

We first prove that $\lVert \tau-\tau_2\rVert_1
\leq 2\, \TV(\tau)$. 
For every $u\in\Z^d$ such that $\tau(u)\neq\tau_2(u)$, necessarily $\lvert \tau(u)-\tau^{\rm rough}_2(u)\rvert\geq\frac{1}{2}$ and hence, 
$$
\left\lvert\tau(u)-\tau_2(u)\right\rvert\leq \left\lvert\tau(u)-\tau^{\rm rough}_2(u)\right\rvert+\frac{1}{2}\leq 2\left\lvert\tau(u)-\tau^{\rm rough}_2(u)\right\rvert.$$
It follows that 
$$ 
\lVert \tau-\tau_2\rVert_1=\sum_{u\in\Z^d}\lvert\tau(u)-\tau_2(u)\rvert\leq 2\sum_{u\in\Z^d}\left\lvert\tau(u)-\tau^{\rm rough}_2(u)\right\rvert.
$$
For every $w\in\Z^d$, by Lemma \ref{lem:diff_rough}, 
$$
\sum_{u\in B_w}\lvert\tau(u)-\tau^{\rm rough}_2(u)\rvert\leq \sum_{\{u,v\}\in E(\Z^d)\cap\binom{B_w}{2}}\lvert\tau(u)-\tau(v)\rvert=\TV\left(\tau;B_w\right),
$$
where $B_w:=Q_2(2\,w)$.
Hence,
\begin{align*} 
\lVert \tau-\tau_2\rVert_1&\leq 2\sum_{u\in\Z^d}\left\lvert\tau(u)-\tau^{\rm rough}_2(u)\right\rvert\\&=2\sum_{w\in\Z^d}\sum_{u\in B_w}\lvert\tau(u)-\tau^{\rm rough}_2(u)\rvert
\leq 2\sum_{w\in\Z^d}\TV\left(\tau;B_w\right)\leq 2\TV(\tau).
\end{align*}

We proceed to show that for any positive integer $N$, it holds that $\lVert \tau_N-\tau_{2N}\rVert_1
\leq (4d+9)N\, \TV(\tau)$.
It is clearly enough to show that for every $w\in \mathbb Z^d$, 
$$
\sum_{u\in Q_{2N}(2N\,w)} \lvert \tau_N(u) -\tau_{2N} (u) \rvert\leq (4d+9)N \TV\left(\tau;Q_{2N}(2N\,w)\right).
$$
Denote $B_w:=Q_2(2\,w)$ once more, and let 
$$A_w:=\left\{v\in B_w\colon \lvert\tau^{\rm rough}_N(N\,v)-\tau_N(N\,v)\rvert>\frac{1}{6}\right\}.$$

For $v\in B_w$ such that $\left\lvert\tau^{\rm rough}_N(N\,v)-\tau^{\rm rough}_{2N}(N\,v)\right\rvert<\frac{1}{3}$, necessarily $\left\lvert\tau_N(N\,v)-\tau_{2N}(N\,v)\right\rvert\leq 1$ and if 
$\left\lvert\tau_N(N\,v)-\tau_{2N}(N\,v)\right\rvert=1$, then $\left\lvert\tau_N(N\,v)-\tau^{\rm rough}_{2N}(N\,v)\right\rvert\geq\frac{1}{2}$ and hence,
$$
\lvert\tau^{\rm rough}_N(N\,v)-\tau_N(N\,v)\rvert\geq \left\lvert\tau_N(N\,v)-\tau^{\rm rough}_{2N}(N\,v)\right\rvert-\left\lvert\tau^{\rm rough}_N(N\,v)-\tau^{\rm rough}_{2N}(N\,v)\right\rvert>\frac{1}{2}-\frac{1}{3}=\frac{1}{6},
$$
i.e., $v\in A_w$;
if $v\in B_w$ such that
$\left\lvert\tau^{\rm rough}_N(N\,v)-\tau^{\rm rough}_{2N}(N\,v)\right\rvert\geq\frac{1}{3}$, then 
$$
\left\lvert\tau_N(N\,v)-\tau_{2N}(N\,v)\right\rvert\leq \left\lvert\tau^{\rm rough}_N(N\,v)-\tau^{\rm rough}_{2N}(N\,v)\right\rvert+1\leq 4\left\lvert\tau^{\rm rough}_N(N\,v)-\tau^{\rm rough}_{2N}(N\,v)\right\rvert.$$
It follows that
\begin{align*}
\frac{1}{N^d}\sum_{u\in Q_{2N}(2N\,w)}\lvert\tau_N(u)-\tau_{2N}(u)\rvert&=
\sum_{v\in B_w}\lvert\tau_N(N\,v)-\tau_{2N}(N\,v)\rvert\\
&\leq \lvert A_w\rvert+
4\sum_{v\in B_w}\lvert\tau^{\rm rough}_N(N\,v)-\tau^{\rm rough}_{2N}(N\,v)\rvert,
\end{align*}
and we are done since by Lemma \ref{lem:Ron}, $$\lvert A_w\rvert\leq \frac{9}{N^{d-1}}\sum_{v\in A_w}\TV\left(\tau; Q_N(N\,v)\right)\leq \frac{9}{N^{d-1}}\TV\left(\tau; Q_{2N}(2N\,w)\right)$$
and by Lemma \ref{lem:diff_rough} and Lemma \ref{lem:rough}, 
\begin{align*}
\sum_{v\in B_w}\lvert\tau^{\rm rough}_N(N\,v)-\tau^{\rm rough}_{2N}(N\,v)\rvert&\leq\sum_{\{u,v\}\in E(\Z^d)\cap\binom{B_w}{2}}\lvert\tau^{\rm rough}_N(N\,u)-\tau^{\rm rough}_N(N\,v)\rvert\\
&\leq \frac{1}{N^{d-1}}\sum_{\{u,v\}\in E(\Z^d)\cap\binom{B_w}{2}}\TV\left(\tau; Q_N(N\,u)\cup Q_N(N\,v)\right)\\
&\leq \frac{d}{N^{d-1}}\TV\left(\tau; Q_{2N}(2N\,w)\right).\qedhere\end{align*}
\end{proof}

\end{document}
